\section{Antecedentes de la empresa}
\subsection{Datos Operativos}
\begin{itemize}
    \item ¿Cuál es el volumen aproximado de material que procesan (toneladas/día)? 30-60 toneladas.
    \item ¿Cuántas sedes tienen? 1 sola.
    \item ¿Cuántos empleados interactúan con sistemas informáticos vs. personal puramente operativo? 3 empleados son los únicos con sistema informático. Operadora de báscula y tickets, contador y administrativo.
\end{itemize}

\subsection{Procesos Críticos}
\begin{itemize}
    \item \textbf{Flujo de información:} Los camiones y chatarra llegan con la operadora de báscula, la cual se asegura del contenido y el tipo de precio que debe ser, se lo comunica a la administradora y cuando todo está en orden, se genera el ticket y la administradora da el dinero para pagarlo. Después de haberlo pagado o de ser necesaria la factura se le dan los datos a la contadora para que se asegure de que las facturas y cuentas sean correctas.
    \item \textbf{Tipos de materiales:} Manejan todo tipo de residuo no peligroso menos plásticos, incluyendo electrónicos. La trazabilidad que manejan con las empresas es mediante manifiestos de SEGAM que se hacen al momento de recoger la chatarra a una empresa y no pueden irse sin uno por lo que únicamente se añade la trazabilidad de llevar el documento hasta el patio de almacenamiento.
\end{itemize}

